
\documentclass{article}
\usepackage{colortbl}
\usepackage{makecell}
\usepackage{multirow}
\usepackage{supertabular}

\begin{document}

\newcounter{utterance}

\centering \large Interaction Transcript for game `cladder', experiment `full\_v1.5\_default', episode 7046 with Qwen3.6{-}27B.
\vspace{24pt}

{ \footnotesize  \setcounter{utterance}{1}
\setlength{\tabcolsep}{0pt}
\begin{supertabular}{c@{$\;$}|p{.15\linewidth}@{}p{.15\linewidth}p{.15\linewidth}p{.15\linewidth}p{.15\linewidth}p{.15\linewidth}}
    \# & \multicolumn{2}{c}{Player} && \multicolumn{2}{c}{Game Master} \\
    \hline

    \theutterance \stepcounter{utterance}  
    & & & \multicolumn{4}{p{0.6\linewidth}}{
        \cellcolor[rgb]{0.9,0.9,0.9}{
            \makecell[{{p{\linewidth}}}]{
                \texttt{\tiny{[P1$\langle$GM]}}
                \texttt{You are an expert in causal inference. The following question is not a typical commonsense query, but rather a meticulously designed question created by a professor specializing in causal inference, intended to assess the students' mastery of the course content.} \\
\texttt{Imagine a self{-}contained, hypothetical world with only the following conditions, and without any unmentioned factors or causal relationships: Jyka has a direct effect on lirg. Hwax has a direct effect on lirg. The overall probability of jyka is 14\%. For those who are not jyka, the probability of lirg is 71\%. For those who are jyka, the probability of lirg is 66\%. Is lirg less likely than not lirg overall?} \\
\texttt{Start your answer with 'Yes' or 'No', followed by additional reasoning or evidence to support your explanation.} \\
            }
        }
    }
     \\ \\

    \theutterance \stepcounter{utterance}  
    & \multicolumn{4}{p{0.6\linewidth}}{
        \cellcolor[rgb]{0.561,0.737,0.561}{
            \makecell[{{p{\linewidth}}}]{
                \texttt{\tiny{[P1$\rangle$GM]}}
                \texttt{No} \\
\\ 
\texttt{To determine whether `lirg` is less likely than `not lirg` overall, we must calculate the total probability of `lirg`, denoted as \$P(L)\$.} \\
\\ 
\texttt{**1. Identify the given probabilities:**} \\
\texttt{*   \$P(J) = 0.14\$ (Probability of `jyka`)} \\
\texttt{*   Therefore, \$P(\textbackslash{}neg J) = 1 {-} 0.14 = 0.86\$ (Probability of `not jyka`)} \\
\texttt{*   \$P(L | \textbackslash{}neg J) = 0.71\$ (Probability of `lirg` given `not jyka`)} \\
\texttt{*   \$P(L | J) = 0.66\$ (Probability of `lirg` given `jyka`)} \\
\\ 
\texttt{**2. Apply the Law of Total Probability:**} \\
\texttt{The overall probability of `lirg` is the sum of the probabilities of `lirg` occurring in each mutually exclusive state of `jyka`:} \\
\\ 
\texttt{\$\$P(L) = P(L | J) \textbackslash{}cdot P(J) + P(L | \textbackslash{}neg J) \textbackslash{}cdot P(\textbackslash{}neg J)\$\$} \\
\\ 
\texttt{**3. Perform the calculation:**} \\
\texttt{\$\$P(L) = (0.66 \textbackslash{}cdot 0.14) + (0.71 \textbackslash{}cdot 0.86)\$\$} \\
\\ 
\texttt{Calculate the first term:} \\
\texttt{\$\$0.66 \textbackslash{}cdot} \\
            }
        }
    }
    & & \\ \\

    \theutterance \stepcounter{utterance}  
    & & & \multicolumn{2}{p{0.3\linewidth}}{
        \cellcolor[rgb]{0.9,0.9,0.9}{
            \makecell[{{p{\linewidth}}}]{
                \texttt{\tiny{[GM$|$GM]}}
                \texttt{No} \\
            }
        }
    }
    & & \\ \\

    \theutterance \stepcounter{utterance}  
    & & & \multicolumn{2}{p{0.3\linewidth}}{
        \cellcolor[rgb]{0.9,0.9,0.9}{
            \makecell[{{p{\linewidth}}}]{
                \texttt{\tiny{[GM$|$GM]}}
                \texttt{no} \\
            }
        }
    }
    & & \\ \\

    \theutterance \stepcounter{utterance}  
    & & & \multicolumn{2}{p{0.3\linewidth}}{
        \cellcolor[rgb]{0.9,0.9,0.9}{
            \makecell[{{p{\linewidth}}}]{
                \texttt{\tiny{[GM$|$GM]}}
                \texttt{game\_result = WIN} \\
            }
        }
    }
    & & \\ \\

\end{supertabular}
}

\end{document}
